\section{The \sys{} Benchmark}
\label{sec:benchmark}

\paragraph{Claim banks.}
Each of four domains is a fictional project followed across sixteen meetings: a clinic booking-system migration, a city flood-barrier programme, a multi-site field-study consortium and a consumer-hardware launch.
A domain's bank holds 24 claims.
Twelve are \emph{critical}: warnings about a risk (e.g.\ a backup that skips a database), conditions that must hold before an action, and dissent from the agreed plan.
Twelve are \emph{incidental}: true but decision-irrelevant details (a room change, a lunch order).
Each claim has four paraphrased mentions, a multiple-choice question whose answer is stated only by that claim, the correct option and two plausible distractors (Appendix~\ref{app:bank}).
A bank also supplies a \emph{majority thread} of routine statements about the plan, filler small talk, a per-session progress line, and eight speakers.

\paragraph{Dialogue construction.}
A dialogue has sixteen sessions and carries six critical and six incidental claims.
Each claim is assigned a frequency $f\in\{1,2,4\}$ by a Latin rotation over dialogue index, so that across every three consecutive dialogues each claim appears at each frequency, and claim content is perfectly balanced across frequency levels.
The session of the claim's last mention, $L$, is drawn uniformly from $\{4,\dots,8\}$ independently of $f$; mentions are spread evenly over $[1,L]$ (\emph{spread} schedule) or placed in the $f$ consecutive sessions ending at $L$ (\emph{clustered} schedule, \S\ref{sec:findings}).
The $k$-th mention uses the $k$-th paraphrase, so single-mention claims are always phrased by the first paraphrase; all claims are therefore equally easy to recognise in isolation.
Every session also contains two majority-thread utterances, three filler utterances and a progress line, rendered as a transcript with speaker names.
We generate 48 dialogues per domain.
A second, naturalistic variant rewrites each session into free-flowing conversation with an LLM while preserving the claims (\S\ref{sec:findings}, \todo{results pending}).

\paragraph{Memory conditions.}
At each session $t$ a memory $M_t$ is produced from $M_{t-1}$ and the transcript of session $t$.
\emph{Rolling} asks the model to rewrite the notes to incorporate the new meeting in at most $B$ words ($B\in\{80,160\}$); outputs are hard-clipped at $B$ words.
\emph{Truncation} keeps the last $B$ words of the transcript, a recency-only control with the same budget.
\emph{Full} gives the reader the whole transcript up to $t$, and \emph{empty} gives no memory, estimating guessability.
The prompts are in Appendix~\ref{app:prompts}.

\paragraph{Probing.}
After sessions $t\in\{4,8,12,16\}$ (full context) and every session (other conditions), the reader answers each claim's question from $M_t$ alone, choosing among the answer, two distractors and ``This was not mentioned in the notes''; option order is randomised with a fixed seed per (dialogue, claim).
A probe is correct ($y=1$) only if the reader picks the answer.
Unless stated otherwise the reader is the same model that writes the memory; \S\ref{sec:findings} reports a fixed external reader.
We analyse only probes after the claim's last mention ($t\ge L$), so every claim has been fully stated by probe time.

\paragraph{Models.}
Qwen2.5-1.5B-Instruct, Qwen2.5-7B-Instruct, Phi-3.5-mini-instruct and Gemma-4-E4B-it, greedy decoding; 7B-scale models are run with 4-bit NF4 weights, and a bf16 comparison on one domain leaves conclusions unchanged (Appendix~\ref{app:quant}).
